


% ---------------------------------------------------

% \subsection{Thin Lens Model (Radiometric)}
% \label{sec:thin_lens}
% Under the thin lens assumption, the lens is a disk of finite radius and infinitesimal width \cite{pharr2023physically}. Therefore, the terms thin lens and aperture are used equivalently in the following.

% \textbf{Defocus}
% All rays from a sensor position $\vbp\in\mathcal{I}$ through the thin lens intersect in a single point on the focal plane~\cite{pharr2023physically}. When $\vbp$ sees a surface that does not lie on the focal plane, the rays from $\vbp$ through the aperture intersect the scene surface in a circle of confusion, leading to a blurry defocus appearance~\cite{pharr2023physically}.

% \textbf{Vignetting}
% The solid angle $\calOmega(\vbp)$ becomes smaller towards the sensor edges. This leads to an intensity falloff to the sensor border and is, among other effects, called vignetting~\cite{pharr2023physically}.
% We can separate this effect from the solid angle by changing the integration domain to the aperture area.
% Let therefore $s'$ be the image distance.
% When integrating over the aperture area $\mathcal{A}$ by using the transformation $\vbomega = (\vba - \vbp)/\lVert \vba - \vbp \rVert$~\cite{pharr2023physically},
% the per-channel radiant flux \cref{eq:per-channel-rad-flux} for a thin lens becomes
% \begin{align}
%     &\Phi_c(\vbu)\label{eq:rad-flux-disent}
%     = \int_{\mathcal{I}_{\vbu}}\int_{\mathcal{A}}\frac{\cos^4\alpha(\vba-\vbp)}{s'^2}\\
%     &\qquad\qquad\qquad L_c(\vbx(\vbp, \vba-\vbp), \vbx(\vbp, \vba-\vbp)-\vba)\odif{\vba}\odif{\vbp}\nonumber.
% \end{align}
% The $\cos^4\alpha(\vba-\vbp)$ term is a well-known physical lens property and describes the vignetting for a thin lens in this formulation~\cite{Goossens_2019}.
% We will utilize the disentangled formulation from \cref{eq:rad-flux-disent} in \cref{sec:vignetting}.







% ---------------------------------------
